\section{Discussion}

\subsection{Ce que le gain signifie}

Le lot fixe ram\`ene la recherche au niveau du self-play : environ 0,4 \`a
0,6\,ms par position contre 0,38\,ms pour le banc self-play \`a lot fixe 8,
l\`a o\`u la recherche en payait 1,4 \`a 2,9. Les trois positions convergent
vers 1\,500 \`a 2\,100 simulations par seconde, alors qu'elles s'\'etalonnaient
de 350 \`a 850 auparavant. Le bot b\'en\'eficie directement de ce facteur, car
l'\'evaluateur occupait 96 \`a 98\,\% du temps.

Le r\'esultat d\'epasse l'estimation de l'\'etude de co\^ut, qui annon\c cait
$\times 3$ \`a $\times 6$ pour l'ensemble du chantier mais pla\c cait la
premi\`ere \'etape, la plomberie, \`a $\times 3$ \`a $\times 5$ au mieux : elle
\'etait juste, et elle est d\'esormais acquise.

\subsection{Le plafond restant}

Le banc \'evaluateur montre que le co\^ut par position continue de baisser
jusqu'au lot 128 : 0,10\,ms contre 0,63\,ms au lot 8, soit un facteur six. Mais
l'arbre d'un seul coup ne fournit que 5 \`a 7 feuilles par vague, et le
padding \'evalue donc des doublons : au remplissage observ\'e de 6,6, environ
17\,\% des rang\'ees d'un lot sont des duplications. Deux directions restent :

\begin{enumerate}[leftmargin=1.4em,itemsep=3pt]
  \item \textbf{Remplir plus.} La divergence a montr\'e qu'on ne remplit pas
  plus sans casser la qualit\'e. Une file d'inf\'erence continue, qui
  accumulerait les requ\^etes pendant le Run, ne forcerait rien et pourrait
  remplir 12 \`a 16 feuilles. Gain estim\'e : $\times 1{,}3$ \`a $\times 1{,}7$
  sur l'\'evaluateur, donc sur la recherche.
  \item \textbf{Descendre le co\^ut par position.} M\^eme \`a forme fixe, le
  GPU reste sous-occup\'e : 0,4\,GFLOP par position pour 3,4\,ms par appel de
  8, soit quelques pour cent de la puissance cr\^ete. Le chemin passe par des
  noyaux mieux adapt\'es au 8$\times$8 (TensorRT, fusion des couches), pas par
  la quantification, qui a empir\'e les choses.
\end{enumerate}

Par rapport \`a Leela Chess Zero, l'\'ecart s'est r\'eduit sans dispara\^itre.
lc0 reste devant sur la taille de lot effective (32 \`a 256 contre 8), le
backend et la file asynchrone. Le pr\'esent chantier a supprim\'e le facteur
qui n'\'etait ni algorithmique ni mat\'eriel : le contrat de forme impos\'e par
ONNX Runtime.

\subsection{Les le\c cons}

\begin{resume}[Quatre le\c cons transf\'erables]
\begin{enumerate}[leftmargin=1.4em,itemsep=2pt]
  \item \textbf{Un contrat implicite peut co\^uter plus cher que
  l'algorithme.} Rien dans le code ne disait que la forme du lot devait \^etre
  stable ; le runtime, lui, la re-planifiait \`a chaque changement.
  \item \textbf{Mesurer hors de l'arbre.} Tant que l'appel est mesur\'e au
  milieu de la recherche, la cause reste noy\'ee. Un banc d'API isol\'e a
  tranch\'e en quelques minutes ce que trois hypoth\`eses se disputaient
  depuis des semaines.
  \item \textbf{Le protocole interleaved n'est pas un luxe.} La d\'erive entre
  sessions est du m\^eme ordre que les gains cherch\'es ; sans alternance, les
  conclusions auraient pu \^etre invers\'ees.
  \item \textbf{La barri\`ere qualit\'e prot\`ege des victoires empoisonn\'ees.}
  Le balayage de divergence offrait 13 \`a 25\,\% de d\'ebit pour 2 points de
  r\'esolution ; seule la comparaison appari\'ee pouvait le montrer.
\end{enumerate}
\end{resume}

\subsection{Les limites}

\begin{itemize}[leftmargin=1.2em,itemsep=2pt]
  \item Une seule machine et une seule carte ; les rapports absolus ne se
  transposent pas telles quelles, mais le m\'ecanisme de re-planification est
  une propri\'et\'e d'ONNX Runtime, pas du GPU.
  \item Le balayage de divergence repose sur 6 observations par position et le
  pr\'efiltre qualit\'e sur 500 puzzles ; une diff\'erence de moins d'un point
  n'aurait pas \'et\'e distingu\'ee du bruit.
  \item Le lot fixe \'evalue des doublons : \`a remplissage 6,6, environ
  17\,\% des rang\'ees sont perdues. C'est le prix de la forme stable, tr\`es
  inf\'erieur au facteur quatre gagn\'e.
  \item Le chemin mono est acc\'el\'er\'e mais reste le fallback ; le
  self-play, qui \'evaluait d\'ej\`a des lots fixes de 128 \`a 512, n'avait
  rien \`a corriger.
  \item Le bot utilise un mod\`ele plus r\'ecent que celui des mesures ; le
  lot fixe ne change pas la recherche, donc la conclusion se transpose, mais
  les d\'ebits absolus diff\`erent.
\end{itemize}

\subsection{La production continue}

Le verdict provisoire, \'etabli dans le rapport de r\'esultats, est le
suivant : la production continue peut viser un remplissage de 12 \`a 16
feuilles et un gain de $\times 1{,}3$ \`a $\times 1{,}7$, pas davantage. Elle
exige de synchroniser les statistiques lues pendant les backups, un co\^ut et
un risque de correction qui m\'eritent un microbenchmark de file avant tout
chantier. La recommandation est d'\'ecrire ce microbenchmark, puis un plan
s\'epar\'e avec les m\^emes crit\`eres de m\'ediane, de p95 et de qualit\'e que
les vagues.
